\documentclass[12 pt, letterpaper]{hmcpset}
\usepackage{hw-preamble}

\name{Jayden Thadani}
\class{MATH 176 --- Algebraic Geometry}
\assignment{Homework assignment \#11}
\duedate{22nd November 2024}

\begin{document}

For this week's problems let $\mathbb{P}^{n}(F)$ denote projective $n$-space over a field $F$. Recall that $S = F[x_{1}, \dots, x_{n}, x_{0}]$ is a \emph{graded ring}, that is, $S = \sum_{e \ge 0} S_{e}$ where the collection $S_{e}$ of homogeneous polynomials of degree $e$ satisfies (i) $S_{0} = F$, (ii) $S_{i} \cap S_{j} = \{ 0 \}$ whenever $i \neq j$, and (iii) $S_{i} S_{j} \subset S_{i + j}$.

\begin{problem}[Problem 1.]
	Let $T \subset S$ be an ideal. Show that the following are equivalent:
	\begin{enumerate}[label = (\roman*)]
		\item Each $f \in T$ can be uniquely expressed as $f = f_{0} + f_{1} + \dots + f_{n}$ where each $f_{e} \in T \cap S_{e}$.
		\item $T = \sum_{e \ge 0} T \cap S_{e}$ and $(T \cap S_{i}) \cap (T \cap S_{j}) \subset \{ 0 \}$ whenever $i \neq j$.
		\item $T \subset \sum_{e \ge 0} T \cap S_{e}$.
	\end{enumerate}
	Any $T$ satisfying either of these is said to be a \emph{homogeneous set}, and we write ${T = \bigoplus_{e \ge 0} T \cap S_{e}}$.
\end{problem}

\begin{solution}
	We will show (i) $\iff$ (ii) $\iff$ (iii).

	Note that each $S_{j}$ is closed under scalar multiplication by elements of $F$ and is also closed under addition (by $S_{0} S_{j} = S_{j}$ and by definition respectively).

	Suppose (i) --- that each  $f \in T$ can be uniquely written as $f = f_{0} + f_{1} + \dots + f_{n}$ with $f_{e} \in S_{e}$. Then clearly $T = \sum_{e \ge 0} T \cap S_{e}$ since each we have written each $f$ as the sum of elements in $T \cap S_{e}$. Further, we can see that $(T \cap S_{i}) \cap (T \cap S_{j}) \subset \{ 0 \}$ when $i \neq j$ since else we could write $0$ in two different ways --- as $0 + \dots + 0$ or as $0 + \dots + 0 + f + \dots - f + 0 + \dots 0$ where $f$ is in both $T \cap S_{i}$ and $T \cap S_{j}$.

	Now suppose (ii). Since $T$ is the sum of its intersections with the $S_{e}$, we can write any $f \in T$ as the sum of $f_{e} \in S_{e}$. It only remains to show that this sum is unique. Suppose that we can write
	\[
		f = f_{0} + f_{1} + \dots + f_{n} = g_{0} + g_{1} + \dots + g_{n}.
	\]
	Then,
	\[
		0 = (f_{0} - g_{0}) + (f_{1} - g_{1}) + \dots + (f_{n} - g_{n})
	\]
	Suppose by way of contradiction that $i$ is the first index for which $f_{i} - g_{i} \neq 0$. Then
	\[
		f_{i} - g_{i} = (g_{i + 1} - f_{i + 1}) + \dots + (g_{n} - f_{n})
	\]
	giving us $S_{i} \cap \sum_{e > i} S_{e} \neq \{ 0 \}$. However, this is false, and thus, all the $f_{i} - g_{i} = 0$ giving us that the sum is unique.

	Again, suppose (ii). Since $T$ is equal to the sum of all the intersections, it is certainly a subset of them. Now suppose (iii). Note that we have $S_{i} \cap S_{j} = \{ 0 \}$ from the definition of $S$. Thus, clearly, $(T \cap S_{i}) \cap (T \cap S_{j}) \subset S_{i} \cap S_{j} \subset \{ 0 \}$ when $i \neq j$. Since $T$ is an ideal, it contains every linear combination of its elements, and thus, contains $\sum_{e \ge 0} T \cap S_{e}$. Thus, by double containment, $T = \sum_{e \ge 0} T \cap S_{e}$.
\end{solution}

\pagebreak

\begin{problem}[Problem 2.]
	Let $\{ T_{\alpha} \mid \alpha \in \mathcal{I} \}$ denote a family of homogeneous sets $T_{\alpha} \subset S$ indexed by some set $\mathcal{I}$.
	\begin{enumerate}
		\item Show that the ideal $T = \left < f \right >$ is a homogeneous ideal for any homogeneous polynomial $f \in S_{e}$.
		\item Show that the union $T = \bigcup_{\alpha \in \mathcal{I}} T_{\alpha}$ of homogeneous ideals $T_{\alpha}$ is also a homogeneous ideal.
		\item Assume $\lvert \mathcal{I} \rvert < \infty$ and write $\mathcal{I} = \{ \alpha_{1}, \dots, \alpha_{r} \}$. Show that the following is also a homogeneous ideal.
			\[
				T = \prod_{\alpha \in \mathcal{I}} T_{\alpha} = \left \{ \sum_{i_{1}, i_{2}, \dots, i_{r}} c_{i_{1}, i_{2}, \dots, i_{r}} f_{\alpha_{1}} f_{\alpha_{2}} \cdots f_{\alpha_{r}} \mid f_{\alpha} \in T_{\alpha} \text{ for each } \alpha \in \mathcal{I} \right \}
			\]
	\end{enumerate}
\end{problem}

\begin{solution}
	\begin{enumerate}
		\item Since $S$ is graded, every polynomial $g \in S$ can be written as the sum of homogeneous components uniquely. Say $g = g_{0} + g_{1}, \dots + g_{m}$ with $g_{i}$ homogeneous of degrees $i$ with $i$ in $1$ through $m$. Thus, $fg$ can be written as the sum of homogeneous components ---
			\[
				fg = fg_{0} + fg_{1} + \dots + fg_{m}
			\]
			with $f g_{i}$ of degree $ei$. That is,
			\[
				\left < f \right > \subset \sum_{e \ge 0} \left < f \right > \cap S_{e}
			\]
			which is sufficient to say that $f$ is homogeneous.
		\item Note that since $T = \bigcup_{\alpha \in \mathcal{I}} T_{\alpha}$, we have $T \cap S_{e} = \bigcup_{\alpha \in \mathcal{I}} T_{\alpha} \cap S_{e}$ by distributing the intersection over the unions. But then we can write
			\begin{equation*}
				\sum_{e \ge 0} T \cap S_{e} = \sum_{e \ge 0} \bigcup_{\alpha \in \mathcal{I}} T_{\alpha} \cap S_{e}.
			\end{equation*}
			It suffices to show
			\[
				\bigcup_{\alpha \in \mathcal{I}} \sum_{e \ge 0} T_{\alpha} \cap S_{e} \subset \sum_{e \ge 0} \bigcup_{\alpha \in \mathcal{I}} T_{\alpha} \cap S_{e}
			\]
			since
			\[
				T = \bigcup_{\alpha \in \mathcal{I}} T_{\alpha} = \bigcup_{\alpha \in \mathcal{I}} \sum_{e \ge 0} T_{\alpha} \cap S_{e}
			\]
			giving us
			\[
				T \subset \sum_{e \ge 0} \bigcup_{\alpha \in \mathcal{I}} T_{\alpha} \cap S_{e} = \sum_{e \ge 0} T \cap S_{e}.
			\]
			But this inclusion just follows since any $f \in \bigcup_{\alpha \in \mathcal{I}} \sum_{e \ge 0} T_{\alpha} \cap S_{e}$ is just an element of $\sum_{e \ge 0} T_{\alpha} \cap S_{e}$ for some $\alpha \in \mathcal{I}$ which can be obtained as an element of $\sum_{e \ge 0} \bigcup_{\alpha \in \mathcal{I}} T_{\alpha} \cap S_{e}$ by picking the same $\alpha$ in the union for each $e$.
		\item Since each $f_{\alpha_{i}} \in T_{\alpha_{i}}$ can be written uniquely as the sum of homogeneous parts, by expanding out, we can write each $f \in T$ as the sum of homogeneous parts. 
	\end{enumerate}
\end{solution}

\pagebreak

\begin{problem}[Problem 3.]
	Given a homogeneous ideal $T \subset S$, define $Z(T) = \{ p \in \mathbb{P}^{n}(F) \mid f(p) = 0 \text{ for all } f \in T \}$.
	\begin{enumerate}
		\item Say that $\{ T_{\alpha} \mid \alpha \in \mathcal{I} \}$ is an arbitrary family of homogeneous ideals. Show
			\[
				Z(T) = \bigcap_{\alpha \in \mathcal{I}} Z(T_{\alpha}) \quad \text{whenever} \quad T = \bigcup_{\alpha \in \mathcal{I}} T_{\alpha}
			\]
		\item Say $\{ T_{\alpha} \mid \alpha \in \mathcal{I} \}$ is a \emph{finite} family of homogeneous ideals. Show
			\[
				Z(T) = \bigcup_{\alpha \in \mathcal{I}} Z(T_{\alpha}) \quad \text{whenever} \quad T = \prod_{\alpha \in \mathcal{I}} T_{\alpha}
			\]
	\end{enumerate}
\end{problem}

\begin{solution}
	\begin{enumerate}
		\item Suppose $p \in Z(T)$. Thus, $f(p) = 0$ for all $f \in T$ and since $T$ is the union of all $T_{\alpha}$, $f(p) = 0$ for all $f \in T_{\alpha}$ for all $\alpha \in \mathcal{I}$. Thus, $p \in Z(T_{\alpha})$ for all $\alpha \in \mathcal{I}$. That is, $p \in \bigcap_{\alpha \in \mathcal{I}} Z(T_{\alpha})$. Thus, $Z(T) \subset \bigcap_{\alpha \in \mathcal{I}} Z(T_{\alpha})$.

			Now suppose $p \in \bigcap_{\alpha \in \mathcal{I}} Z(T_{\alpha})$. Since $f(p) = 0$ for all $f \in T_{\alpha}$ for all $\alpha \in \mathcal{I}$ and the elements of $T$ are exactly those functions, we also have $f(p) = 0$ for all $f \in T$. Thus, $p \in Z(T)$. Now, by double containment,
			\[
				Z(T) = \bigcap_{\alpha \in \mathcal{I}} Z(T_{\alpha}).
			\]
		\item Let $\mathcal{I} = \{ \alpha_{1}, \dots, \alpha_{\lvert \mathcal{I} \rvert} \}$.

			Suppose $p \in Z(T)$. Suppose by way of contradiction that $f \notin Z(T_{\alpha})$ for any $\alpha \in \mathcal{I}$. Notice that $Z(T)$ contains all products $f = f_{\alpha_{1}} \cdots f_{\alpha_{\lvert \mathcal{I} \rvert}}$ and we can choose $f_{\alpha_{i}}$ all not to vanish at $p$ (this is our supposition). Then we have $f \in T$ with
			\[
				f(p) = f_{\alpha_{1}}(p) \cdots f_{\alpha_{\lvert \mathcal{I} \rvert}}(p) \neq 0.
			\]
			which contradicts $p \in Z(T)$. Thus our supposition is false and $p \in Z(T_{\alpha})$ for some $\alpha \in \mathcal{I}$. That is, $p \in \bigcup_{\alpha \in \mathcal{I}} Z(T_{\alpha})$.

			Suppose $p \in \bigcup_{\alpha \in \mathcal{I}} Z(T_{\alpha})$. Thus, all $f$ in some $T_{\alpha}$ vanish on $p$. Since every $h \in T$ is a product of $h = f g$, we have $h(p) = f(p) g(p) = 0$ for all $h \in T$. That is, $p \in Z(T)$. Thus, by double containment,
			\[
				Z(T) = \bigcup_{\alpha \in \mathcal{I}} Z(T_{\alpha}).
			\]
	\end{enumerate}
\end{solution}

\pagebreak

\begin{problem}[Problem 4.]
	We say $U \subset \mathbb{P}^{n}(F)$ is \emph{Zariski open} if there is some homogeneous $T$ such that $U = \mathbb{P}^{n}(F) \setminus Z(T)$.
	\begin{enumerate}
		\item Show that both $\O$ and $\mathbb{P}^{n}(F)$ are Zariski open.
		\item Say $\{ U_{\alpha} \mid \alpha \in \mathcal{I} \}$ is an arbitrary family of Zariski open sets. Show $\bigcup_{\alpha \in \mathcal{I}} U_{\alpha}$ is also Zariski open.
		\item Say $\{ U_{\alpha} \mid \alpha \in \mathcal{I} \}$ is a \emph{finite} family of Zariski open sets. Show $\bigcap_{\alpha \in \mathcal{I}} U_{\alpha}$ is also Zariski open.
	\end{enumerate}
	Denoting the family of all Zariski open sets by $\{ U_{\alpha} \mid \alpha \in \mathcal{I} \}$, we say $(\mathbb{P}^{n}(F), \{ U_{\alpha} \mid \alpha \in \mathcal{I} \})$ is the \emph{Zariski topology on projective space}
\end{problem}

\begin{solution}
	\begin{enumerate}
		\item Note that $\O = Z(S)$ since there are no points on which all polynomials in  $F[x_{0}, x_{1}, \dots, x_{n}]$ vanish, and $\mathbb{P}^{n}(F) = Z(\left < 0 \right >)$ since the zero polynomial vanishes everywhere. But  
			\[
				\O =  \mathbb{P}^{n}(F) \setminus \mathbb{P}^{n}(F) = \mathbb{P}^{n}(F) \setminus Z(\left < 0 \right >)
			\]
			and
			\[
				\mathbb{P}^{n}(F) = \mathbb{P}^{n}(F) \setminus \O = \mathbb{P}^{n}(F) \setminus Z(S).
			\]
			Thus, both sets are open.
		\item We can write $U_{\alpha} = \mathbb{P}^{n}(F) \setminus Z(T_{\alpha})$ for each $\alpha \in \mathcal{I}$. Note that by De Morgan's laws
			\[
				\bigcup_{\alpha \in \mathcal{I}} U_{\alpha} = \bigcup_{\alpha \in \mathcal{I}} \mathbb{P}^{n}(F) \setminus Z(T_{\alpha}) = \mathbb{P}^{n}(F) \setminus \bigcap_{\alpha \in \mathcal{I}} Z(T_{\alpha}).
			\]
			By the previous exercise, we can write $\bigcap_{\alpha \in \mathcal{I}} Z(T_{\alpha}) = Z(T)$ for $T = \bigcup_{\alpha \in \mathcal{I}} T_{\alpha}$. Thus, $\bigcup_{\alpha \in \mathcal{I}} U_{\alpha} = \mathbb{P}^{n}(F) \setminus Z(T)$ is Zariski open.
		\item We can write $U_{\alpha} = \mathbb{P}^{n}(F) \setminus Z(T_{\alpha})$ for each $\alpha \in \mathcal{I}$. Note that by De Morgan's laws
			\[
				\bigcap_{\alpha \in \mathcal{I}} U_{\alpha} = \bigcap_{\alpha \in \mathcal{I}} \mathbb{P}^{n}(F) \setminus Z(T_{\alpha}) = \mathbb{P}^{n}(F) \setminus \bigcup_{\alpha \in \mathcal{I}} Z(T_{\alpha}).
			\]
			By the previous exercise, we can write $\bigcup_{\alpha \in \mathcal{I}} Z(T_{\alpha}) = Z(T)$ for $T = \prod_{\alpha \in \mathcal{I}} T_{\alpha}$ when $\mathcal{I}$ is a finite indexing set. Thus, $\bigcap_{\alpha \in \mathcal{I}} U_{\alpha} = \mathbb{P}^{n}(F) \setminus Z(T)$ is Zariski open.
	\end{enumerate}
\end{solution}

\end{document}
